\section{Hạn chế và hướng phát triển}\label{sec:limits}

\subsection{Hạn chế}
\begin{itemize}
\item \textbf{Dữ liệu là dạng bảng, không phải chuỗi tự nhiên.} Biểu diễn A tạo ``chuỗi'' bằng cách xếp 11 đặc trưng theo một thứ
  tự tuỳ ý; RNN/CNN không có lợi thế nào so với MLP và thua mô hình cây. Biểu diễn B là chuỗi thật (theo thời gian) nhưng chỉ có ích
  cho địa chỉ được dùng lại --- 90\% dòng test là địa chỉ xuất hiện lần đầu, chuỗi chỉ dài 1.
\item \textbf{Lớp \code{white} được lấy mẫu} (khoảng 1.000 dòng/ngày trong dữ liệu gốc), nên ``lịch sử'' của một địa chỉ white bị
  thưa hơn thực tế, và đặc trưng ``số lần đã xuất hiện'' có thể phân biệt hai lớp tốt hơn so với khi triển khai thật trên toàn bộ
  chuỗi khối. Kết quả của biểu diễn B có thể lạc quan so với thực tế.
\item \textbf{Nhãn ransomware không đầy đủ}: chỉ các địa chỉ đã được các nghiên cứu trước xác định; nhiều dòng \code{white} có thể là
  ransomware chưa được phát hiện, làm giảm precision quan sát được.
\item \textbf{Hai cách chia khác nhau}: biểu diễn A dùng chia ngẫu nhiên (để so sánh với Bài 1) --- 54,7\% dòng ransomware trong test
  có địa chỉ đã xuất hiện trong train, nên kết quả A có rò rỉ nhẹ qua địa chỉ. Bảng~\ref{tab:group_other} đã kiểm chứng: trên cách chia
  theo địa chỉ, A cho kết quả gần như giống hệt, nên kết luận không thay đổi.
\item \textbf{Ước lượng phương sai còn hạn chế}: chỉ 3 seed cho các cấu hình chính (1 seed cho các biến thể còn lại), và chỉ một cách
  chia dữ liệu (không cross-validation). Khác biệt $<$0,01 giữa các cấu hình không nên coi là có ý nghĩa.
\item \textbf{Early stopping theo \code{val\_us}} (31,6 nghìn dòng, chỉ $\approx$6,6 nghìn dòng ransomware) có nhiễu; nhiều cấu hình A
  cùng seed dừng ở cùng epoch (42) vì dùng chung thứ tự batch, cho thấy chọn epoch phụ thuộc phần nào vào dao động ngẫu nhiên.
\item \textbf{Siêu tham số chưa tối ưu toàn diện}: thay đổi từng yếu tố một so với cấu hình gốc (không tìm kiếm lưới/Bayes);
  không thử recurrent dropout, weight decay, lịch learning rate. Thời gian/epoch trong bảng chính bị ảnh hưởng do chạy song song
  (đã đo lại riêng ở Bảng~\ref{tab:bench}).
\end{itemize}

\subsection{Hướng phát triển}
\begin{itemize}
\item \textbf{Kết hợp hai thế mạnh}: mô hình cây cho đặc trưng của dòng hiện tại và RNN/CNN cho lịch sử --- ví dụ dùng embedding
  lịch sử do CNN-B học được làm đặc trưng đầu vào cho XGBoost, hoặc ensemble xác suất của hai mô hình (XGBoost tốt hơn ở địa chỉ
  mới, CNN-B tốt hơn ở địa chỉ có lịch sử dài --- Bảng~\ref{tab:by_len}).
\item \textbf{Mô hình đồ thị}: ransomware chuyển tiền qua nhiều địa chỉ; cần dữ liệu giao dịch gốc (cạnh giữa các địa chỉ) để dùng
  Graph Neural Network --- giúp cả với địa chỉ dùng một lần như Cerber/Locky.
\item Transformer trên lịch sử địa chỉ (self-attention với mã hoá khoảng cách thời gian); mô hình FT-Transformer [8] cho biểu diễn A.
\item Tối ưu siêu tham số tự động (Optuna), nhiều seed hơn, cross-validation theo địa chỉ, và hiệu chỉnh xác suất (temperature scaling)
  trước khi tìm ngưỡng/trọng số lớp.
\item Đánh giá theo thời gian (train trên các năm trước, test trên năm sau) để mô phỏng triển khai thực tế và kiểm tra khả năng phát hiện
  họ ransomware mới.
\end{itemize}
